\documentclass[11pt,a4paper]{article}

\usepackage[utf8]{inputenc}
\usepackage[T1]{fontenc}
\usepackage[english]{babel}
\usepackage{ctex}
\usepackage[margin=2.5cm]{geometry}
\usepackage{parskip}
\usepackage{amsmath}
\usepackage{amssymb}
\usepackage{bm}
\usepackage[numbers,sort&compress]{natbib}
\renewcommand*{\bibfont}{\footnotesize}
\bibliographystyle{gbt7714-numerical}
\usepackage{xcolor}
\usepackage{hyperref}
\hypersetup{
  colorlinks=true,
  linkcolor=blue,
  urlcolor=cyan,
  pdftitle={供者来源游离 DNA
与多组学标志物在实体器官移植排斥监测中的应用价值与未来方向},
  pdfsubject={Systematic Literature Review},
}

\title{供者来源游离 DNA
与多组学标志物在实体器官移植排斥监测中的应用价值与未来方向}
\author{}
\date{2026-06-03}

\providecommand{\tightlist}{%
  \setlength{\itemsep}{0pt}\setlength{\parskip}{0pt}}

\begin{document}
\maketitle
\tableofcontents
\newpage

\begin{abstract}
供者来源游离 DNA（donor-derived cell-free DNA,
dd-cfDNA）已经成为近五年实体器官移植无创监测中最接近临床常规化的分子工具之一。它对移植物细胞损伤高度敏感，可纵向重复采样，并能在肾、心、肺、肝等器官中与排斥、微血管炎症、供者特异性抗体（donor-specific antibody,
DSA）和病理损伤形成关联。现有证据的共同结论并不是``dd-cfDNA
等于排斥''，而是 dd-cfDNA
构成了一个跨器官的损伤轴：低位且稳定的结果有助于支持免疫静息或低活检收益，高位、持续升高或相对基线突增则提示需要寻找免疫性、感染性、缺血性或机械性损伤来源。局限也同样明确：dd-cfDNA
反映的是``损伤释放''，并不天然区分抗体介导排斥、T
细胞介导排斥、缺血再灌注损伤、感染、药物毒性、胆汁淤积或近期手术损伤。未来 3-5
年，更可行也更有增量的路线，并不是把 dd-cfDNA
简单作为活检替代品，而是将其与 DSA、尿液外泌体
mRNA、转录组、蛋白组、病毒负荷、免疫抑制暴露、术前机器灌注标志物和可解释 AI
模型整合为动态风险框架，用于判断何时活检、何时减免疫、何时强化治疗，以及何时避免过度诊断。
\end{abstract}

\section{引言}\label{ux5f15ux8a00}

传统移植监测依赖血清肌酐、蛋白尿、肝酶、心功能、肺功能、药物浓度、DSA
和组织活检。问题在于这些证据各自只覆盖风险链条的一段：功能指标往往滞后且不特异，DSA
能提示供受者免疫识别和未来风险但不能直接量化当前组织损伤，活检虽然仍是诊断和分型的核心，却存在侵入性、采样误差、病灶异质性和读片差异。dd-cfDNA
的临床吸引力来自它把移植物细胞死亡转化为外周血可重复测量的信号，并且半衰期短，理论上更接近实时损伤状态。ADMIRAL
等纵向研究推动了肾移植 dd-cfDNA 从研究工具走向临床监测；Trifecta
研究进一步把 dd-cfDNA
与分子病理联系起来，提示它更像``损伤强度计''而非单一排斥标签\cite{clinicaloutcomesfrom2021, thetrifectastudy2022, distinctmolecularprocesses2023}。

这种定位决定了综述的重点应从``能否诊断排斥''转向``在什么临床上下文中改变决策''。若患者低危、功能稳定、DSA
阴性且 dd-cfDNA 长期低位，检测的主要价值是提高排除活动性损伤的把握；若 DSA
新发、微血管炎症风险高或既往 ABMR，dd-cfDNA
升高更可能提高活检收益；若同时存在病毒复制、胆道并发症、肺部感染、机械循环支持或近期缺血再灌注，dd-cfDNA
阳性则首先提示需要更宽的病因鉴别。Banff
最小侵入诊断工作组也强调，血液和尿液标志物目前更适合作为组织学和临床判断的辅助，而不是脱离场景的独立诊断工具\cite{diagnosticpotentialof2022, donorderivedcell20251}。

近年证据显示，dd-cfDNA
正在心、肺、肝和多器官移植中迅速扩展，但不同器官的基线、释放动力学、取样频率和干扰因素差异明显\cite{clinicalutilityof2024, circulatingcellfree2024, donorderivedcell2023}。肾移植的证据最丰富，因而适合讨论证据边界；心移植和肺移植更强调替代频繁侵入性活检和排除高风险事件；肝移植则受到肝细胞体量、再生、胆汁淤积和围手术期释放的影响；多器官移植进一步挑战``一个血浆比例对应一个移植物''的直觉。因此，dd-cfDNA
的临床价值需要放在器官特异性场景中解释，不能用同一阈值横跨所有移植类型。

\section{肾移植：证据最丰富，也最能暴露边界}\label{ux80beux79fbux690dux8bc1ux636eux6700ux4e30ux5bccux4e5fux6700ux80fdux66b4ux9732ux8fb9ux754c}

肾移植是 dd-cfDNA 临床验证最充分的领域，也是最能暴露解释边界的场景。ADMIRAL
研究显示，dd-cfDNA 纵向升高可与临床和亚临床排斥、dnDSA
形成以及后续 eGFR 下降相关；低位结果则支持移植物静息和低活动性损伤概率\cite{clinicaloutcomesfrom2021}。这一证据使
dd-cfDNA
适合进入常规随访的风险分层：它可在肌酐尚未明显变化时提示潜在损伤，也可在功能稳定但免疫风险升高的患者中帮助决定是否提前活检。重要的是，这种价值更多来自连续监测和临床整合，而不是某一个跨平台固定阈值。

ABMR 是肾移植中 dd-cfDNA
表现最稳定的应用场景。Trifecta 研究显示，dd-cfDNA
与分子 AMR、微血管炎症和组织损伤关系紧密；即便 DSA 阴性，分子或组织学 AMR
也可伴随 dd-cfDNA 升高，提示 dd-cfDNA
能捕捉 DSA 未覆盖的损伤后果\cite{thetrifectastudy2022, antibodymediatedrejection2022}。临床研究还显示
dd-cfDNA 与 dnDSA、特别是高 MFI 或 class II dnDSA
联合时，能更好识别晚期同种免疫损伤和 ABMR/MVI 风险\cite{donorderivedcell2022, combiningdonorderived20221, donorderivedcell20242}。在
AT1R 抗体等非 HLA 抗体场景中，dd-cfDNA
与 Banff 微血管炎症评分的关系也提示它可用于判断抗体存在是否已经转化为组织损伤\cite{utilityofdonor2021}。

随机试验的出现使证据边界更清楚。针对已有 dnDSA
的肾移植受者，dd-cfDNA
指导活检可缩短 AMR 诊断时间，说明它不仅有诊断相关性，也可能改变诊断路径\cite{donorderivedcell20241}。但是，这并不等于可以把所有
dd-cfDNA 升高都直接当作 AMR
治疗。TCMR、急性肾损伤、BK 病毒相关肾病、缺血再灌注、钙调神经磷酸酶抑制剂毒性、复发性肾病和近期手术损伤均可能升高
dd-cfDNA；BKPyVAN 还可能同时伴随尿液趋化因子升高，使其与 TCMR
的非侵入性模式重叠\cite{levelsofdonor2022, diagnosticpotentialof2022}。因此，肾移植中的合理策略是把
dd-cfDNA 用作``是否存在实质损伤''和``是否提高活检收益''的证据，而不是单独完成病因分型。

fraction 与 absolute quantity
是肾移植解释中最重要的技术边界之一。fraction 反映供者 cfDNA
占总 cfDNA 的比例，便于跨样本理解，但会受受者总 cfDNA
背景影响；手术、炎症、感染、肥胖、白细胞释放或其他系统性损伤可使分母升高，从而把真实供者释放``稀释''成较低比例。absolute
quantity 或 copies/mL
试图直接量化供者来源拷贝数，可在高总 cfDNA 背景下补充分数值的假阴性风险。Trifecta
相关分析和后续队列均支持同时使用 fraction 与 quantity
或连续复合评分，以提高对活动性排斥的识别能力\cite{usingboththe2021, combiningdonorderived2022, improvingdiagnosticperformance2025}。不过绝对量也受采血、保存、提取、UMI
计数和平台校准影响；不同方法之间百分比可高度相关，绝对拷贝数仍可能存在比例偏倚\cite{comparisonofmethods2023, dropletdigitalpcr2023}。因此，报告系统最好同时给出
fraction、absolute quantity、总 cfDNA 背景、个体化基线和相对变化，而不是只给一个``阳性/阴性''。

\section{心、肺、肝与多器官移植：从可测到可解释}\label{ux5fc3ux80baux809dux4e0eux591aux5668ux5b98ux79fbux690dux4eceux53efux6d4bux5230ux53efux89e3ux91ca}

心移植长期依赖心内膜活检，但活检频繁、侵入性强，且对局灶病变和读片一致性存在限制。dd-cfDNA
在心移植中的直接吸引力，是减少低收益监测活检并更早识别 ABMR
或移植物损伤。Trifecta-Heart 研究显示，血浆 dd-cfDNA
不仅与分子排斥通路相关，也与损伤和巨噬细胞浸润相关，说明心肌细胞损伤本身可形成强烈炎症背景\cite{comparingplasmadonor2024}。前瞻性多测试研究进一步显示，分子
ABMR 与 dd-cfDNA、DSA 的关系强于传统组织学 ABMR；当组织学、分子诊断、DSA
和 dd-cfDNA 不一致时，dd-cfDNA
可帮助判断哪一组证据更接近当前损伤状态\cite{definingtherelationships2025}。

心移植的临床解释更依赖``前测概率''。在 DSA
阴性、心功能稳定的患者中，低 dd-cfDNA
更适合排除高风险事件；在 DSA 阳性、左室功能下降或疑似 AMR
患者中，同样的 dd-cfDNA
升高会显著提高活检和治疗评估的价值\cite{redefiningcardiacantibody2024, donorderivedcell20253}。单中心实践提示，以
dd-cfDNA 为核心的现代监测路径可能减少心内膜活检数量且不增加短期不良结局，但也可能识别更多 AMR
相关信号，需要避免把每一次低度升高都转化为治疗\cite{amodernheart2024, initiationofnoninvasive2021}。心衰、感染、缺血损伤、机械循环支持和人群差异会改变
cfDNA 背景；因此，心移植中的 dd-cfDNA
应和 GEP、DSA、超声/磁共振、肌钙蛋白、临床症状和活检结果共同解释\cite{racialdifferencesin2024, europeansocietyfor2024}。

肺移植的优势和困难同时更突出。肺移植受者感染、原发移植物功能障碍、急性细胞排斥、AMR、CLAD
和气道并发症频繁发生，组织采样负担高且病灶异质。临床验证研究显示，dd-cfDNA
对急性排斥具有较高阴性预测价值，也能反映感染和 CLAD/NRAD
等多种同种移植物损伤\cite{clinicalvalidationof2022, cellfreedna2022}。这使其非常适合做``低风险排除''和``何时进一步支气管镜/活检''的触发器，却不适合单独区分排斥与感染。对于移植两年以上患者，相对变化可能比固定绝对阈值更能识别急性肺移植物功能障碍，提示肺移植尤其需要个体基线和参考变化值\cite{relativechangein2023}。系统综述和单中心研究也支持
dd-cfDNA 在肺移植排斥监测中的诊断潜力，但异质性较高，未来试验必须把感染、CLAD
和治疗反应分别建模\cite{circulatingdonorderived2023, applicationofplasma2023, whycellfree2022}。

肝移植则不能简单套用肾或心的阈值。肝脏体量大、细胞更新和再生能力强，围手术期 dd-cfDNA
峰值高，胆汁淤积、缺血再灌注、血管/胆道并发症、感染和药物性肝损伤均可能导致升高。成人肝移植研究显示，术后
dd-cfDNA 与转氨酶和组织损伤相关，D-HOPE
灌注液中的线粒体 cfDNA 还与早期移植物功能和严重并发症相关\cite{circulatingcellfree2024}。儿科和前瞻性队列提示
dd-cfDNA 可辅助识别急性排斥，但胆汁淤积等混杂因素会降低阳性预测值；miRNA
等肝损伤/免疫调节标志物可能提高病因解释力\cite{preliminaryclinicalexperience2021, donorderivedcell20252, advancesandchallenges2025}。因此，肝移植 dd-cfDNA
更应作为``移植物完整性和损伤负荷''指标，结合肝酶模式、胆道影像、病毒学、药物暴露、miRNA
和必要活检判断，而不是单独充当排斥诊断。

多器官移植进一步提出来源归属问题。当心肾、肝肾、胰肾或心肺联合移植时，血浆 dd-cfDNA
可能来自多个供体器官，比例值受到器官体量、基线释放和并发症的共同影响。心脏多器官移植研究显示，心肝和心肺受者
dd-cfDNA 水平和波动明显高于心肾，且可长期超过心移植常用阈值而无心脏排斥\cite{clinicalutilityof2024}。MOTR
研究也提示，不同多器官组合的基线和器官功能异常关联不同，LKT 等组合的 dd-cfDNA
水平不能按肾单器官阈值解释\cite{donorderivedcell2026}。未来若要在多器官受者中做精确决策，需要供体特异基因型、组织来源甲基化/片段组学、器官特异蛋白或转录标志物共同判断损伤来自哪个移植物。

\section{多组学整合：让 dd-cfDNA
从``损伤读数''变成``病因推断''}\label{ux591aux7ec4ux5b66ux6574ux5408ux8ba9-dd-cfdna-ux4eceux635fux4f24ux8bfbux6570ux53d8ux6210ux75c5ux56e0ux63a8ux65ad}

dd-cfDNA
最大的临床短板是病因特异性不足，多组学的价值也正在这里体现出来。一个可执行的整合框架可以把证据分为四层：第一层是损伤负荷，包括
dd-cfDNA fraction、absolute quantity、总 cfDNA
和趋势；第二层是免疫风险，包括 DSA、非 HLA
抗体、HLA/eplet mismatch、既往排斥和免疫不依从；第三层是炎症和病因线索，包括尿液外泌体
mRNA、尿液 CXCL9/CXCL10、外周血转录组、蛋白组、miRNA、病毒负荷和感染证据；第四层是组织锚定，包括病理、MMDx
或其他分子病理分类。只有当多个层面的证据方向一致时，dd-cfDNA
才应推动强化免疫治疗；若证据分叉，合理动作往往是复测、寻找感染或进行病理确认。

DSA 与 dd-cfDNA 的整合最接近临床落地。DSA
提示供受者之间存在免疫识别，dd-cfDNA
提示这种风险是否已经转化为当前损伤。class II dnDSA、高 MFI DSA、DQ
抗体或 DSA 伴移植物功能异常时，dd-cfDNA
升高显著提高 ABMR/MVI 的前测概率；反之，DSA 阳性但 dd-cfDNA
长期低位可能提示暂不需要立即治疗，而应加强监测或选择性活检\cite{donorderivedcell2022, combiningdonorderived20221, perspectivefordonor2024}。但是
DSA 阴性并不排除 ABMR，非 HLA
抗体和分子 AMR 均可能释放 dd-cfDNA；因此联合模型不应把 DSA
作为绝对门槛，而应把 DSA、dd-cfDNA 和病理/分子表型视为互补证据\cite{antibodymediatedrejection2022, utilityofdonor2021}。

尿液和外泌体组学适合补足肾移植的局部病因线索。尿液外泌体、mRNA/miRNA、趋化因子和尿液 dd-cfDNA
更接近泌尿系统局部微环境，可提供肾小管炎症、病毒性损伤和免疫细胞募集信息\cite{liquidbiopsyfor2023, molecularimmunemonitoring2023}。已有研究显示，尿液
dd-cfDNA 浓度可能更明显关联 BK
病毒相关肾病，而血浆 fraction 更突出 ABMR 和微血管炎症；尿液趋化因子在 BKPyVAN
与 TCMR 中又可能重叠\cite{usingbothplasma2022, levelsofdonor2022}。这说明尿液信号不能简单作为血浆
dd-cfDNA 的替代，而应作为病因鉴别的一部分：血浆 dd-cfDNA
回答``移植物是否在受损''，尿液外泌体/趋化因子帮助回答``肾内炎症、病毒复制或小管损伤是否参与''。

外周血转录组、蛋白组和分子病理提供另一个维度。免疫静息基因表达签名与 dd-cfDNA
信号不完全相关，二者联合可改善对肾移植静息与排斥的区分\cite{clinicalvalidationof2021}。真实世界研究也提醒，单个商业转录组或
dd-cfDNA 检测未必能在未选择人群中准确诊断所有排斥，但在 DSA
阳性且移植物功能异常的 for-cause 场景中，dd-cfDNA
与 DSA 组合对 AMR/MVI 更有价值\cite{bloodgeneexpression2024}。MMDx
等分子病理进一步显示，AMR、TCMR、近期实质损伤和萎缩纤维化具有不同转录表型，dd-cfDNA
可作为校准这些组织状态的血液损伤读数\cite{themolecularphenotype2023, distinctmolecularprocesses2023}。蛋白组、代谢组和 cfDNA
甲基化/片段组学未来可能把``损伤来源''和``损伤通路''拆得更细，但当前仍需要标准化和外部验证\cite{areviewof2024, donorderivedcell20251}。

病毒负荷和免疫抑制暴露必须进入同一个决策模型。对肾移植而言，BK
病毒载量、尿液趋化因子、dd-cfDNA 和肌酐变化共同决定是减免疫、活检还是抗病毒/支持治疗；对心肺移植而言，CMV、呼吸道病毒、细菌/真菌感染和
dd-cfDNA 同时升高时，贸然强化免疫抑制可能造成更大风险。TTV
等病毒负荷和细胞功能测定可作为过度免疫抑制风险的反向轴，与 dd-cfDNA
代表的拒斥/损伤轴共同服务精准免疫抑制\cite{predictingoverimmunosuppression2024}。临床上最有用的模型不是简单给出``排斥概率''，而是同时估计排斥风险、感染风险、免疫抑制过度风险和活检收益。

机器灌注标志物可成为移植前层面的先验。扩展标准供体器官、DCD
器官或高缺血风险移植物在植入前已经释放代谢、炎症、线粒体损伤和内皮损伤信号；若这些灌注液指标与植入后
dd-cfDNA
动态连接，就可能区分``可预期的保存损伤''与``异常免疫损伤''。肝移植 D-HOPE
灌注液 mt-cfDNA 与移植后功能和严重并发症的关联，为这种术前-术后连续建模提供了直接线索\cite{circulatingcellfree2024}。在肾、肝、心肺等高风险供体场景中，机器灌注读数可作为
AI 风险模型的先验，使术后早期 dd-cfDNA
峰值不被过度解释为排斥，从而减少不必要活检和过早治疗。

\section{讨论：临床应用场景与解释边界}\label{ux8ba8ux8bbaux4e34ux5e8aux5e94ux7528ux573aux666fux4e0eux89e3ux91caux8fb9ux754c}

临床应用可以先落在五类场景。第一是稳定低风险随访：若 dd-cfDNA
长期低位、DSA 阴性、功能稳定且尿液/转录组无炎症信号，协议活检的边际收益可能下降，复查间隔可更个体化。第二是高风险筛查：免疫不依从、DSA
新发、既往 ABMR、HLA class II eplet mismatch
高、边缘供体器官或感染后免疫调节患者，dd-cfDNA
趋势可帮助决定活检时点和随访密度\cite{donorderivedcell20232, elevationofdonor2024}。第三是
for-cause 诊断：当肌酐、蛋白尿、肝酶、肺功能或心功能异常时，dd-cfDNA
可以提高对实质损伤的判断，但病因仍需结合 DSA、病毒、影像和组织学。第四是治疗反应监测：ABMR
或 TCMR 治疗后，dd-cfDNA
下降可能比单一功能指标更快反映组织损伤缓解；若功能改善但 dd-cfDNA
持续升高，则应警惕残余微血管炎症、持续感染或其他隐匿损伤\cite{associationofblood2025, donorderivedcell20244}。第五是免疫抑制精准化：dd-cfDNA
代表低免疫抑制/排斥风险的一端，病毒负荷、TTV、机会感染、肿瘤和细胞功能测定代表过度免疫抑制的一端，二者共同决定减药、维持或强化治疗。

过度诊断风险贯穿所有场景。dd-cfDNA
检测越敏感，越容易发现亚临床、短暂或非排斥性损伤；如果临床路径把每一次升高都转化为活检或抗排斥治疗，就可能增加医疗负担、患者焦虑、感染风险和药物毒性。心移植中，GEP
和 dd-cfDNA 联合监测曾出现检测升高导致心内膜活检增加但阳性收益有限的情况\cite{initiationofnoninvasive2021}。肝移植中，胆汁淤积可造成
dd-cfDNA 升高并降低阳性预测值\cite{donorderivedcell20252}。肺移植中，感染和 CLAD
可与排斥共享损伤信号\cite{clinicalvalidationof2022}。因此，检测前应明确 context
of use：是为了排除活动性损伤、提高活检收益、监测治疗反应，还是支持免疫抑制调整。没有预设动作阈值的检测，会把``信息增加''误变成``临床收益''。

AI 风险模型的价值在于整合这些互相冲突的证据，但模型必须可解释。输入层应包含器官类型、术后时间、供体质量和机器灌注先验、dd-cfDNA
fraction 与 quantity、历史基线和相对变化、DSA/eplet mismatch、尿液/血液转录组、蛋白组或 miRNA、病毒负荷、药物浓度、感染证据和近期干预。输出层不应只有黑箱风险分，而应给出分解后的风险贡献：例如``ABMR/MVI
风险升高，主要由 DSA、dd-cfDNA quantity 和微血管损伤相关证据驱动''，或``损伤信号升高但感染证据强，建议先复测和排除感染后再决定活检''。MMDx
等分子分类已经展示了机器学习解释组织状态的可能性，但临床 AI
模型还需要校准、外部验证、决策曲线分析和前瞻性影响评价\cite{themolecularphenotype2023, improvingdiagnosticperformance2025}。

临床决策试验应从诊断准确性转向患者和系统结局。理想设计不是简单比较 dd-cfDNA
阳性与活检诊断，而是随机或阶梯楔形比较``常规监测''与``dd-cfDNA
多组学决策算法''。入组应按器官和风险分层，包括肾移植 dnDSA/既往 ABMR
人群、心移植低危监测和 DSA 阳性人群、肺移植感染/CLAD
高风险人群、肝移植高缺血或胆道风险人群以及多器官移植。终点应包括活检数量、活检阳性率、诊断延误、过度治疗、感染和肿瘤事件、患者焦虑、成本、eGFR/FEV1/LVEF/肝功能轨迹、移植物存活和死亡。只有当算法减少无效活检和误治，同时不增加漏诊或长期损伤时，dd-cfDNA
才真正从``优秀标志物''变成``有效临床策略''。

\section{展望：未来 3-5
年研究方向}\label{ux5c55ux671bux672aux6765-3-5-ux5e74ux7814ux7a76ux65b9ux5411}

\begin{enumerate}
\def\labelenumi{\arabic{enumi}.}
\tightlist
\item
  建立器官特异、时间特异的动态阈值。早期术后、稳定期、感染期、DSA
  新发期、治疗后和多器官移植不应使用同一解释框架；肺移植尤其需要把相对变化和参考变化值纳入常规解释。
\item
  同时报告 fraction、absolute quantity、总 cfDNA 背景和个体化变化率。比例值容易受受者
  cfDNA 背景影响，绝对量也受平台和前分析因素影响；二者结合比任何单一阈值更接近真实损伤判断。
\item
  开展 dd-cfDNA + DSA + 尿液外泌体/趋化因子 + 转录组/蛋白组 +
  分子病理的前瞻性决策试验。终点不应只看诊断准确性，还应看活检减少、误治减少、排斥延误减少、感染事件和长期功能。
\item
  把机器灌注标志物纳入早期损伤解释。对
  DCD、扩展标准供体和高冷缺血器官，术前灌注液读数可作为 dd-cfDNA
  术后曲线的先验，帮助区分保存损伤、早期功能不全和异常免疫损伤。
\item
  发展可解释 AI
  风险模型。模型应输出``可能病因组合''、``证据冲突点''和``建议下一步证据''，而不是黑箱风险分。临床可用性取决于它能否告诉医生该活检、复查、调整免疫抑制还是寻找感染。
\item
  明确过度诊断的停止规则。连续低位、短暂孤立升高、感染证据充分或术前灌注先验可解释的早期峰值，都应有相应的复测和观察路径，避免把检测敏感性误转化为不必要治疗。
\end{enumerate}

\section{结论}\label{ux7ed3ux8bba}

dd-cfDNA
已经改变移植随访的语言：医生不再只能等待器官功能变差，而能更早看到移植物细胞损伤的分子信号。但
dd-cfDNA
不是活检的简单替代品，也不是排斥的单一同义词。肾移植证据支持其用于损伤监测、ABMR/MVI
风险分层和 dnDSA 场景下的活检触发，但也清楚显示 TCMR、BK
病毒、急性肾损伤和总 cfDNA 背景会限制特异性；心、肺、肝和多器官移植则各有基线、动力学和混杂因素。它更合适的位置，是作为多证据动态监测系统中的核心损伤轴，与免疫风险、感染风险、组织/尿液组学、病毒负荷、免疫抑制暴露和移植前器官质量共同构成临床决策模型。未来
3-5
年，能进入常规路径的方案，应当是那些经过前瞻性决策试验证明能够减少不必要活检、避免过度治疗、降低漏诊风险并改善长期结局的整合型监测策略。

\nocite{*}
\bibliography{ddcfDNA-multiomics-rejection-monitoring_参考文献.bib}

\end{document}
